% GENERATED FILE. Edit canonical lessons, then run npm run generate:books.
% canonical-source-hash: fnv1a64:db78e598f0d9a8c4
% canonical-lessons: PA-R166-run-again, PA-R166-wander-again, PA-R166-stop-again

\chapter{Return to Running, Wandering, and Stopping}
\label{ch:pa-three-familiar-actions-r4}
\hlchaptermodality{166}

\begin{chapteropening}
\textbf{By the end of this chapter:} \emph{I can retrieve three already-taught Punjabi action verbs after a long gap, one at a time, without claiming scored A1 writing or complete mock credit.}

Recall running, wandering, and stopping from short familiar cues without introducing new Punjabi.
\end{chapteropening}

\section[\texorpdfstring{\pa{ਭੱਜਣਾ} (bhajjṇā)}{ਭੱਜਣਾ (bhajjṇā)}]{An old word for running}

\label{lesson:PA-R166-run-again}



\begin{quote}
Hear \textbf{to run}. Recall \textbf{\pa{ਭੱਜਣਾ}} (\emph{bhajjṇā}).
\end{quote}

\subsection*{Your turn}
Hear \textbf{\pa{ਭੱਜਣਾ}}. Recall \textbf{to run}. Then hear \textbf{to run} again. Say \textbf{\pa{ਭੱਜਣਾ}} before the model. Only this old word is needed.

\begin{tcolorbox}[breakable,colback=teal!4,colframe=teal!35!black,title={Before you move on}]
One last cue: \textbf{to run}. Recall \textbf{\pa{ਭੱਜਣਾ}}. Listen and retry once if it does not come back.
\end{tcolorbox}
\section[\texorpdfstring{\pa{ਘੁੰਮਣਾ} (ghummṇā)}{ਘੁੰਮਣਾ (ghummṇā)}]{An old word for wandering}

\label{lesson:PA-R166-wander-again}



\begin{quote}
Hear \textbf{to wander, to go around}. Recall \textbf{\pa{ਘੁੰਮਣਾ}} (\emph{ghummṇā}).
\end{quote}

\subsection*{Your turn}
Hear \textbf{\pa{ਘੁੰਮਣਾ}}. Recall \textbf{to wander, to go around}. Then hear the English cue again. Say \textbf{\pa{ਘੁੰਮਣਾ}} before the model. Only this old word is needed.

\begin{tcolorbox}[breakable,colback=teal!4,colframe=teal!35!black,title={Before you move on}]
One last cue: \textbf{to wander, to go around}. Recall \textbf{\pa{ਘੁੰਮਣਾ}}. Listen and retry once if it does not come back.
\end{tcolorbox}
\section[\texorpdfstring{\pa{ਰੁਕਣਾ} (rukṇā)}{ਰੁਕਣਾ (rukṇā)}]{An old word for stopping}

\label{lesson:PA-R166-stop-again}



\begin{quote}
Hear \textbf{to stop}. Recall \textbf{\pa{ਰੁਕਣਾ}} (\emph{rukṇā}).
\end{quote}

\subsection*{Your turn}
Pause after each familiar meaning, then hear its model:

\begin{itemize}
\raggedright
  \item \textbf{to wander, to go around} — \textbf{\pa{ਘੁੰਮਣਾ}} (\emph{ghummṇā}).
  \item \textbf{to stop} — \textbf{\pa{ਰੁਕਣਾ}} (\emph{rukṇā}).
  \item \textbf{to run} — \textbf{\pa{ਭੱਜਣਾ}} (\emph{bhajjṇā}).
\end{itemize}

Only the order of these old words has changed.

\begin{tcolorbox}[breakable,colback=teal!4,colframe=teal!35!black,title={Before you move on}]
Hear \textbf{to stop}, \textbf{to run}, then \textbf{to wander}. Pause after each cue and recall its Punjabi word before the model: \textbf{\pa{ਰੁਕਣਾ}}, \textbf{\pa{ਭੱਜਣਾ}}, \textbf{\pa{ਘੁੰਮਣਾ}}. Revisit only the old word you missed.
\end{tcolorbox}
